% Guide to the current state of the TFG (DerStandard: replies, stance and user lifetime)
% Build: pdflatex tfg-guide.tex (twice, for the table of contents)
\documentclass[11pt,a4paper]{article}
\usepackage[utf8]{inputenc}
\usepackage[T1]{fontenc}
\usepackage{lmodern}
\usepackage{amsmath,amssymb}
\usepackage[margin=2.4cm]{geometry}
\usepackage{booktabs,tabularx,longtable,array}
\usepackage{enumitem}
\usepackage{xcolor}
\usepackage[hidelinks]{hyperref}
\usepackage{microtype}

\setlist{itemsep=2pt,topsep=4pt}
\setlength{\parindent}{0pt}
\setlength{\parskip}{5pt}
\renewcommand{\arraystretch}{1.15}

\usepackage{url}
\urlstyle{tt}
\newcommand{\file}[1]{\path{#1}}
\newcommand{\todo}[1]{\textcolor{red!70!black}{[#1]}}
\newcommand{\readref}[1]{\textcolor{blue!50!black}{\textbf{Read:}} #1}
\newcommand{\status}[1]{\textsc{\small #1}}
\newcolumntype{Y}{>{\raggedright\arraybackslash}X}

\title{\textbf{A guide through my TFG}\\[2pt]\large Who replies to whom, are they opposed, and how long do they stay?\\ Reply structure $\times$ vote-based stance $\times$ user lifetime on DerStandard.at}
\author{Working guide, state of 2026-11-02}
\date{}

\begin{document}
\maketitle
\tableofcontents

\section*{How to use this guide}
\addcontentsline{toc}{section}{How to use this guide}
This is a map, not a textbook. It says \emph{what} was done, \emph{why}, \emph{where it lives in the repository} and \emph{what to read} to understand each piece. It deliberately contains no statistical theory: wherever a method appears, a \readref{} line points to a source. Citations like [\ref{ref:mh}] refer to the reading list in Section~\ref{sec:reading}.

\textbf{Suggested order.} (1) Sections~\ref{sec:context}--\ref{sec:hyp} give the story. (2) Section~\ref{sec:choices} lists every decision that someone can challenge in a defence. (3) Sections~\ref{sec:process}--\ref{sec:methods} explain how the work was done and what each method answers. (4) Sections~\ref{sec:results}--\ref{sec:conclusions} give numbers and the claims they support. (5) Section~\ref{sec:selftest} is a self-test; if you cannot answer a question, go back to the section it names.

\textbf{Status tags.} The repository tags each decision as \status{[me]}, \status{[AI, reviewed]} or \status{[AI, not yet reviewed]}. The first analysis session (2026-10-06) was largely carried out with an AI assistant (see \file{notes/method/AI-USE.md}). The aim of this guide is that every item below becomes something you can explain from understanding.

\textbf{Caveat on the references.} The bibliographic details in Section~\ref{sec:reading} were written from memory and are not checked online. Verify each DOI/page range before it enters the written thesis. PDFs already in \file{notes/papers/} are marked as such.

% ------------------------------------------------------------------
\section{Context and aim}\label{sec:context}
\subsection{The project in one paragraph}
The TFG is supervised by Andreas Kaltenbrunner (UPF) and is built on the DerStandard.at dataset created by his group (Fraxanet, G\'omez, Kaltenbrunner, Pellert): ten years (2013--2022) of threaded comments and signed (up/down) votes from an Austrian newspaper forum. The first assigned task (meeting 2026-09-30) was: \emph{analyse who replies to whom and whether their positions are opposed, complemented with an analysis of users' lifetime.} The second meeting (2026-10-19) approved the method and results and asked for follow-ups; the third meeting is 2026-11-02.

\subsection{Timeline}
\begin{tabularx}{\textwidth}{@{}l Y@{}}
\toprule
Date & What happened \\
\midrule
2026-09-30 & Meeting 1: ideation. FAULTANA future work discarded (already tested without interesting results). Task set: replies $\times$ opposition $\times$ lifetime. Biweekly meetings. \\
2026-10-06 & Dataset downloaded (28\,GB, MD5-verified), plan written, scripts 00--07 run, results note produced (largely AI-executed). \\
2026-10-16/19 & Null model approved; meeting 2: method and results approved. Follow-ups: positive votes, understand H2 and the methods, top 10\% longest-living users, study signed networks. \\
2026-11-02 & Scripts 08--09 run (positive votes; longest-living users). Meeting 3. \\
\bottomrule
\end{tabularx}

\subsection{Where this sits in the literature}
\begin{itemize}
  \item \textbf{Signed networks and structural balance} (Heider; Cartwright \& Harary; Leskovec et al.). Edges are positive or negative; the question is how the sign structure organises the group. \readref{[\ref{ref:heider}], [\ref{ref:cartwright}], [\ref{ref:leskovec}], and the signed-networks chapter of [\ref{ref:easley}].}
  \item \textbf{FAULTANA} (Fraxanet et al., 2024): detects the antagonism/alignment fault line in a signed vote network of this same forum, giving the two factions used in H3. Local PDF: \file{Fraxanet2024_faultana-antagonism-alignment.pdf} plus its SI. \readref{[\ref{ref:faultana}].}
  \item \textbf{The dataset paper} (Fraxanet et al., 2026): what each file is and its known inconsistencies. Local PDF \file{Fraxanet2025_derstandard-dataset.pdf}. \readref{[\ref{ref:dataset}].}
  \item \textbf{Prior group work on replies and votes}: Lan Ceqin (2026), reciprocity of replies and votes and new-user shares (local PDF \file{LanCeqin2026_reciprocity-derstandard-thesis.pdf}). This TFG's novelty claim is the \emph{cross} of reply structure $\times$ stance $\times$ lifetime, which neither FAULTANA nor that thesis studied.
  \item \textbf{User lifecycle and churn in online communities}: Danescu-Niculescu-Mizil et al.\ (2013), Cheng et al.\ (2015), Kraut \& Resnick (2012). \readref{[\ref{ref:dnm}], [\ref{ref:cheng}], [\ref{ref:kraut}].}
  \item \textbf{Discussion-thread structure}: G\'omez et al.\ (2008, 2013), from the supervisor's group. \readref{[\ref{ref:gomez08}], [\ref{ref:gomez13}].}
\end{itemize}

% ------------------------------------------------------------------
\section{The data}\label{sec:data}
\subsection{What the dataset contains}
DerStandard.at, 2013-01-01 to 2022-12-31 (data collected 2021-06 to 2022-12). Tabular files are TSV with \texttt{NA} for missing values, hashed 16-hex user and comment IDs and unix timestamps.

\begin{tabularx}{\textwidth}{@{}l r Y@{}}
\toprule
Table & Size (raw) & Used for \\
\midrule
Comments & 75.6\,M rows & Author, parent, thread depth, article, timestamp. Reply edges. \\
Votes & 412\,M rows & Voter, comment author, sign ($\pm1$), month only. Stance. \\
Users & 248\,k users & First/last comment and vote. Lifetime. \\
Articles & 587\,k & Topic. Heterogeneity. \\
Summary files & --- & \file{faction_users} (FAULTANA factions, 14.8\,k users: 9\,278 in faction 1, 5\,539 in faction 2); \file{inconsistent_comment_metadata}. \\
Embeddings & 508\,GB & \textbf{Not downloaded} (see choice C1). \\
\bottomrule
\end{tabularx}

\subsection{Properties that shape every analysis}
\begin{enumerate}
  \item \textbf{There is no comment text.} Stance must be inferred from votes.
  \item \textbf{Votes have no timestamp}, only the month of the file. We cannot tell whether a downvote preceded the reply; all stance results are associations.
  \item \textbf{Cleaning}: 259\,521 exact duplicate comments (0.34\%); 27.1\,M votes with a deleted voter; reply edges lost to inconsistent metadata (0.6\%), deleted authors (14.9\%) and self-replies (1.4\%). Final: 75.4\,M unique comments, 383.9\,M usable votes, \textbf{44\,350\,672 reply edges}. Documented in \file{notes/dataset/data-treatment.md}.
  \item \textbf{The FAULTANA factions are survivors by construction}: only users who voted every year 2014--2021 were labelled.
  \item \textbf{Volume change in 2022}: comments and votes per month fall by about 45\% from February 2022 while active users stay flat. Unexplained; to ask Andreas/Emma.
\end{enumerate}
\readref{[\ref{ref:dataset}] for the files and inconsistencies; the local \file{Codebook.md} for field meanings.}

% ------------------------------------------------------------------
\section{Research questions and hypotheses}\label{sec:hyp}
\textbf{Main question.} Who replies to whom, are they in opposed positions, and how is this related to how long users stay in the forum?

Because there is no text, ``opposed position'' means \emph{vote-based disapproval} at three nested levels, from local to global.

\begin{tabularx}{\textwidth}{@{}l l Y Y@{}}
\toprule
 & Level & Statement & Test \\
\midrule
H1 & Comment & Voters who replied to comment $c$ downvote $c$ more often than voters who did not reply. & Within-comment comparison; Mantel--Haenszel pooled odds ratio over 8.1\,M comments. \\
H2 & User pair & A replies to users toward whom A's overall vote balance is more negative than toward a co-present null target. & Beta-posterior edge sign, real vs null, bootstrap clustered by replier. \\
H3 & Faction & The share of cross-faction replies exceeds random mixing and the null; cross-faction replies are more hostile. & Shares vs two baselines; stance by faction relation. \\
H4 & Lifetime & New commenters who receive more downvotes and hostile replies in their first 30 days leave sooner. & Landmark Cox proportional hazards, stratified by cohort. \\
\bottomrule
\end{tabularx}

\textbf{Follow-up questions (meeting 2).} (i) Do the results also hold for \emph{positive} votes? (ii) What distinguishes the longest-living users (top 10\%)? Both have been run (Section~\ref{sec:results}). Still open: signed networks as background reading, and the optional ideas listed in Section~\ref{sec:next}.

\textbf{Important framing.} H1--H3 are associations on observational data; H4 is a hazard association. None is a causal claim. \readref{[\ref{ref:hernan}] for the vocabulary of causal vs associational claims.}

% ------------------------------------------------------------------
\section{Methodological choices and why}\label{sec:choices}
Each row is a decision that can be challenged. ``Where'' gives the repository location.

\begin{longtable}{@{}p{0.04\textwidth} p{0.19\textwidth} p{0.33\textwidth} p{0.28\textwidth}@{}}
\toprule
\# & Choice & Reason & Where / reading \\
\midrule
\endhead
C1 & Do not download the 508\,GB embeddings & Embeddings give semantic similarity, not agreement; the signed vote network measures opposition directly at full scale. & \file{notes/dataset/NOTES.md} \\
C2 & Stance from votes at three levels & Text is absent; three levels answer different questions (this comment, this person, this group) and can disagree. & \file{PLAN.md}, \file{EXPLICACION.md} \S3 \\
C3 & Reply = a message whose \texttt{parent\_id} is B's comment and whose author is A & Only the thread structure certifies a reply; mentions in free text are ignored (accepted loss). & \file{data-treatment.md}, D15 \\
C4 & Remove duplicates, deleted voters, self-replies, inconsistent records & Duplicates double-count edges and votes and break independence. & \file{data-treatment.md}, D7 \\
C5 & \textbf{Co-present null model}: replace B by the author of a uniformly random \emph{earlier} comment in the same article & Keeps article, time and exposure fixed; only the choice of target changes. Without it ``31\% negative relations'' has no reference. & \file{02_null_pairs.py}; \file{NOTES.md} 4 \\
C6 & Stratify by comment (Mantel--Haenszel) in H1 & Compares voters on the \emph{same} comment, so a controversial comment does not inflate the effect. \readref{[\ref{ref:mh}], [\ref{ref:rbg}], [\ref{ref:agresti}].} & \file{04_h1_comment_level.py} \\
C7 & Beta posterior for the pair relation, prior $(\alpha,\beta)=(1,2)$, edge negative if $E[p]<0.4$, positive if $>0.6$ & Inherited from FAULTANA so that relations are comparable. The prior is mildly negative; the multiverse tests alternatives. \readref{[\ref{ref:faultana}], [\ref{ref:bda}].} & \file{05_h2_h3_dyad_faction.py} \\
C8 & Exclude A's vote on the specific comment in H2 & Otherwise H2 re-counts H1. & \file{03_dyads.py} \\
C9 & Per-edge metrics (negative-edge share, mean $E[p]$) instead of the pooled downvote share & The pooled share is dominated by a few mega-dyads with thousands of mutual upvotes (median 1--2 votes per pair, mean $\sim$110). & \file{NOTES.md} 7 \\
C10 & Bootstrap clustered by replier (Poisson weights) & Edges of the same replier are not independent. \readref{[\ref{ref:efron}].} & \file{05_h2_h3_*.py} \\
C11 & Lifetime = last $-$ first activity (comments and votes); right-censored if last activity $\ge$ 2022-07-01; left-censored if first activity $<$ 2013-07-01 & The window ends at 2022-12; users still active may not have left. Left-censored users may predate the data. \readref{[\ref{ref:km}], [\ref{ref:kleinbaum}].} & \file{06_lifetime.py} \\
C12 & \textbf{Landmark design} at day 30 in H4 & Exposures are measured in the first 30 days; counting survival from day 0 would credit users with ``surviving'' while exposures accumulate (immortal-time bias). \readref{[\ref{ref:suissa}], [\ref{ref:anderson}].} & \file{h4.py}, \file{06_lifetime.py} \\
C13 & Cox model stratified by entry cohort & Cohorts differ in baseline hazard (newer cohorts leave faster); stratifying avoids assuming one baseline. \readref{[\ref{ref:cox}], [\ref{ref:kleinbaum}].} & \file{06_lifetime.py} \\
C14 & Shrink early shares toward the \emph{population} rate (pseudo-count $M=2$) & A first version used a Beta(1,1) prior, putting users with no replies at 0.5, above the real base rate. ``No data'' looked hostile and the HR changed sign. Fixed. & \file{NOTES.md} 8 \\
C15 & Multiverse / specification curves for H2 (24 specs) and H4 (18 specs) & Shows how much the headline depends on arbitrary choices. \readref{[\ref{ref:steegen}], [\ref{ref:simonsohn}] (both local PDFs).} & \file{07_multiverse.py} \\
C16 & Report effect sizes and CIs, not p-values & With $N$ in the millions everything is ``significant''. \readref{[\ref{ref:wasserstein}], [\ref{ref:lin}].} & \file{RESULTS.md} \\
C17 & ``Long-lived'' = top 10\% of lifetime \emph{within entry cohort} (preferred) & Top 10\% overall mostly selects early cohorts, which had more time. Also reported: top 10\% overall, $>2$ y, $>3$ y. \status{[AI, not yet reviewed]} & \file{09_top_decile_lifetime.py} \\
C18 & Positive edge = $E[p]>0.6$, mirror of the negative edge & Same prior, thresholds and null; makes positive and negative results comparable. \status{[AI, not yet reviewed]} & \file{08_positive_votes.py} \\
\bottomrule
\end{longtable}

% ------------------------------------------------------------------
\section{Process and reproducibility}\label{sec:process}
\subsection{Repository layout}
One experiment = one folder with the same name in \file{notes/experiments/} (documentation and results) and in \file{code/experiments/} (scripts). The project map is \file{INDEX.md}. The notes folder is an Obsidian vault. Raw and derived data (about 42\,GB) live in \file{data/derstandard/}, outside git.

\subsection{Pipeline}
\begin{tabularx}{\textwidth}{@{}p{0.22\textwidth} p{0.27\textwidth} Y r@{}}
\toprule
Script & Step & Output & Time \\
\midrule
\file{00_to_parquet} & TSV $\to$ Parquet, typed IDs, \texttt{NA}$\to$NULL, votes cleaned & \file{comments, votes, votes_clean, users, articles, factions}.parquet & $\sim$6.5\,min \\
\file{01_reply_edges} & Reply edges with tenure of A and B & \file{reply_edges}.parquet & 20\,s \\
\file{02_null_pairs} & Co-present null target per reply (hash-based, fixed seed) & \file{null_edges}.parquet & 1\,min \\
\file{03_dyads} & Directed vote totals for real and null pairs; per-edge stance & \file{dyads}, \file{edges_stance}.parquet & 1\,min \\
\file{04_h1_comment_level} & H1: Mantel--Haenszel, bootstrap, year/depth/topic & CSV + figure & 35\,s \\
\file{05_h2_h3_dyad_faction} & H2, H3, faction projection, reply-back & CSVs + figure & 30\,s \\
\file{06_lifetime} & Lifetime, Kaplan--Meier, tenure table, H4 Cox & CSVs + figures & 1--2\,min \\
\file{07_multiverse} & Specification curves for H2 and H4 & CSVs + figure & 3\,min \\
\file{08_positive_votes} & Positive-vote mirror of H1--H4 & \file{08_*} CSVs + figure & 3.5\,min \\
\file{09_top_decile_lifetime} & Longest-living users & \file{09_*} CSVs + figure & $<$\,10\,s \\
\bottomrule
\end{tabularx}
\smallskip
Full re-run from raw data takes roughly 15--20 minutes on the project machine (32\,GB RAM). Every step logs row counts and timings to \file{results/run.log}. Helpers: \file{common.py} (paths, censoring constants, logging), \file{h4.py} (the H4 landmark feature builder, used by 06, 07, 08, 09), \file{plotstyle.py}.

\subsection{Tools}
DuckDB for out-of-core SQL over Parquet (the 400\,M-row joins do not fit in pandas), then pandas/NumPy/SciPy/statsmodels for statistics, \texttt{lifelines} for survival models, matplotlib for figures; \texttt{uv} manages the environment. \readref{[\ref{ref:duckdb}], [\ref{ref:lifelines}].}

\subsection{Documentation habits}
Every methodological decision goes into \file{NOTES.md} of the experiment with a status tag; cleaning decisions go into \file{notes/dataset/data-treatment.md}; every number in a result note is traceable to a CSV in \file{results/}. Documents for understanding: \file{EXPLICACION.md} (long), \file{EXPLICACION-RESUMEN.md} (7-minute version), \file{PREGUNTAS.md} (practice questions).

% ------------------------------------------------------------------
\section{Methods in plain words}\label{sec:methods}
For each method: the question it answers in this project, one-sentence intuition, a common pitfall, and what to read.

\begin{description}[leftmargin=0pt,style=nextline]
\item[Odds ratio and Mantel--Haenszel (H1).]
Question: do repliers downvote more than other voters \emph{on the same comment}? The odds ratio compares odds of a downvote between the two groups; Mantel--Haenszel pools one $2\times2$ table per comment into a single ratio, so differences between comments cancel out. Pitfall: an OR of 4.4 is \emph{not} ``4.4 times more likely''; quote the percentages (53\% vs 30\%) beside it. For a single table, the upvote OR is exactly the inverse of the downvote OR. \readref{[\ref{ref:mh}], [\ref{ref:rbg}] (the variance used), [\ref{ref:agresti}].}

\item[Null model (H2, H3).]
Question: what would we see if A had replied to somebody who was merely present? Intuition: a counterfactual draw that keeps article, time and exposure and changes only who is chosen. Pitfall: the null target can coincide with the real one (4\% of cases), which is legitimate. \readref{Section on permutation/randomisation baselines in [\ref{ref:wasserman_net}]; see also how FAULTANA builds its baselines [\ref{ref:faultana}].}

\item[Beta posterior for a pair (H2).]
Question: is A \emph{in general} for or against B? The probability that A votes positively on B is given a Beta prior and updated by counting A's votes on B; the posterior mean $E[p]$ and its variance define the edge sign. Pitfall: with one vote, the prior dominates (hence the ``$\ge5$ votes'' variant). \readref{[\ref{ref:bda}] (beta--binomial model); [\ref{ref:faultana}] for the prior and thresholds used.}

\item[Bootstrap clustered by replier.]
Question: how uncertain is the difference real vs null when a replier contributes many correlated edges? Intuition: resample users, not edges. Poisson(1) weights are a fast equivalent. \readref{[\ref{ref:efron}].}

\item[Kaplan--Meier and censoring (lifetime).]
Question: what fraction of users is still active after $t$ years, counting users who have not yet left? \readref{[\ref{ref:km}], [\ref{ref:kleinbaum}].}

\item[Cox proportional hazards, landmark, stratification (H4).]
Question: does early hostility change the instantaneous risk of leaving? The hazard ratio is per standard deviation of the exposure. The landmark at day 30 avoids immortal-time bias; stratifying by cohort lets each cohort have its own baseline. The proportional-hazards assumption is checked with Schoenfeld residuals. Pitfall: HR $>1$ means leaving \emph{sooner}. \readref{[\ref{ref:cox}], [\ref{ref:schoenfeld}], [\ref{ref:grambsch}], [\ref{ref:suissa}], [\ref{ref:anderson}].}

\item[Shrinkage of user-level shares (H4).]
Question: how to compute ``share of hostile replies'' for a user with 0, 1 or 2 replies? Intuition: pull small samples toward the population rate with a pseudo-count. See choice C14 for the bug this fixed. \readref{Empirical-Bayes sections of [\ref{ref:bda}].}

\item[Multiverse / specification curve.]
Question: do conclusions survive reasonable alternative analytic choices? Re-run with every defensible combination and plot the spread. \readref{[\ref{ref:steegen}], [\ref{ref:simonsohn}] (local PDFs, with the latter's SI).}

\item[Logistic regression with cohort fixed effects (longest-living users).]
Question: does early negativity predict being in the long-lived group, once the cohort is controlled for? \readref{[\ref{ref:agresti}].}

\item[Signed networks and structural balance (background for H3).]
Intuition: friends of friends should be friends; enemies of enemies friends; a balanced network splits into two opposed camps, which is the idea behind the two factions. \readref{[\ref{ref:heider}], [\ref{ref:cartwright}], [\ref{ref:leskovec}], [\ref{ref:easley}].}
\end{description}

% ------------------------------------------------------------------
\section{Results}\label{sec:results}
All numbers below are from \file{RESULTS.md} (sections 3--4 and the 2026-11-02 update in section 8) and the CSV tables in \file{results/}.

\subsection{Objective 1: who replies to whom, and are they opposed?}
\begin{tabularx}{\textwidth}{@{}l Y@{}}
\toprule
 & Result \\
\midrule
H1 & On the same comment, repliers downvote 53.3\% vs 30.1\% for other voters; $\mathrm{OR}_{MH}=4.43$ [4.42, 4.44] over 8\,117\,960 comments. Stable by year (4.1--4.7) and topic (3.5--4.9). Escalates with thread depth: OR 3.5 at top level $\to$ 5.7 $\to$ 6.9 $\to$ 7.8 $\to$ 8.2 $\to$ 9.0 at depth $\ge5$. \\
H2 & Negative relations (pairs with $\ge5$ votes, $E[p]<0.4$): 31.4\% for the replied-to user vs 18.1\% for the null, a difference of +13.3 pp [12.6, 14.0] (votes on the replied comment excluded). Strict FAULTANA edges: 9.4\% vs 3.8\%. Across 24 specifications: +11.7 to +20.6 pp, positive in 24/24. \\
H3 & Among 11.3\,M replies with both users labelled: cross-faction 26.0\% observed vs 23.2\% null vs 46.8\% under random mixing. Exposure is strongly segregated; within it, a small excess of out-group targets, coming from the majority faction (16.7\% vs 13.7\%). Cross-faction replies are far more hostile: replier downvoted the parent in 64.6\% vs 37.0\% for same-faction replies. \\
Reply-back & P(B replies to A) is 31.8\% if A downvoted B's comment, 29.8\% if A did not vote, 16.6\% if A upvoted. \\
Tenure & Veterans (3+ y) downvote newcomers ($<30$ d) in 68\% of replies, and other veterans in 48\%. \\
\bottomrule
\end{tabularx}

\subsection{Objective 2: user lifetime and H4}
\begin{tabularx}{\textwidth}{@{}l Y@{}}
\toprule
 & Result \\
\midrule
Descriptives & 247\,863 users (211\,108 commenters). 26.7\% active for less than a day, 37.9\% less than 30 days, 20.3\% wrote one comment; 30.3\% right-censored, 16.6\% left-censored. \\
Cohorts & Median lifetime of commenters (Kaplan--Meier): 1.36 y (2014), 0.81 (2016), 0.41 (2018), 0.39 (2020). One-year survival 53\% $\to$ 41\%. \\
H4 & 98\,670 new commenters (59\% leave before 2022-07). Hazard ratio of leaving per SD of first-30-day exposure: downvote share received 1.11 [1.10, 1.12]; hostile replies received 1.06 [1.05, 1.07]; hostile replies written 1.04; got any reply 0.95; log early comments 1.23. Proportional hazards hold; CI excludes 1 in 18/18 specifications for the three hostility measures. \\
\bottomrule
\end{tabularx}
Note: very active new accounts (more than 30 comments in the first month) leave much faster (71\% events); probably bursts, bans or sockpuppets. Not separated yet.

\subsection{Follow-ups from meeting 2 (run on 2026-11-02, \status{[AI, not yet reviewed]})}
\textbf{Positive votes.} Mirror image of the negative results: repliers upvote the comment they answer less (46.7\% vs 69.9\%; $\mathrm{OR}_{up}=0.23=1/4.43$, not independent evidence). Replied-to users have a lower positive-edge share than the null (68.6\% vs 81.9\%, $-13.3$ pp; $-20.6$ to $-11.7$ pp across 24 specifications, negative in all). Same-faction replies upvote the parent in 63.0\% of cases, cross-faction in 35.4\%. For H4, early upvote share received is protective (HR 0.87 per SD), but friendly replies add nothing once hostility is controlled for (HR $\approx1.00$). It is the \emph{absence of hostility}, not the presence of friendliness, that matters.

\textbf{Longest-living users.} Population: the H4 landmark population (98\,670 users). Top 10\% within cohort, versus the rest, in the first 30 days: downvote share received 20.4\% vs 30.0\%; hostile replies received 11.8\% vs 16.9\%; hostile replies written 12.2\% vs 17.1\%; same early activity (median 2 comments) and same chance of getting a reply (62.6\% vs 62.8\%). Odds of being long-lived per SD: downvote share received 0.87 [0.85, 0.89], hostile replies received 0.94, written 0.96. Inside the group early negativity barely relates to how long they stay (Spearman $\rho$ between 0.03 and 0.07). Over their whole life, at every tenure they downvote the parent less (47--50\% vs 52--54\%) and receive fewer hostile replies. The ``>2 y'', ``>3 y'' and ``top 10\% overall'' definitions give the same direction.

% ------------------------------------------------------------------
\section{Conclusions, limits and what you can claim}\label{sec:conclusions}
\subsection{Conclusions}
\begin{enumerate}
  \item \textbf{Replies are mostly an antagonistic channel.} People reply to comments they downvote and to users they are generally against, more than to other users present in the same discussion (H1, H2; positive-vote results are the mirror image).
  \item \textbf{Conversations escalate}: antagonism grows with thread depth and being downvoted roughly doubles the chance of an answer.
  \item \textbf{The fault line shapes who meets whom}: factions mostly discuss separately; when they meet, replies are hostile.
  \item \textbf{Most newcomers leave quickly, and later cohorts leave faster.}
  \item \textbf{Early hostility is associated with leaving sooner}, robust across specifications. Long-lived users were treated and behaved less negatively from the start.
\end{enumerate}

\subsection{What you can and cannot claim}
\begin{tabularx}{\textwidth}{@{}Y Y@{}}
\toprule
Can claim & Cannot claim \\
\midrule
Repliers disapprove (by vote) of what they answer much more than other voters on the same comment. & That disagreement \emph{causes} replying: votes have no timestamp, so the order is unknown. \\
Early hostility is associated with a higher hazard of leaving, robustly to specification. & That hostility \emph{causes} leaving (reverse causation, unobserved bans), or that veterans do ``gatekeeping'' (compatible, untested). \\
Factions mostly meet within their own faction. & That faction hostility is independent evidence: the partition was built from the same votes (partly circular). \\
Long-lived users start less negative. & That the top 10\% are the top 10\% of the \emph{process}: lifetime is right-censored and the group is selected on survival. \\
\bottomrule
\end{tabularx}

\subsection{Limitations}
\begin{itemize}
  \item ``Position'' is vote-based disapproval, not measured ideology; whether a vote is a valid proxy of opposition is an open question.
  \item Votes carry no timestamp; everything in Objective 1 is an association.
  \item Faction labels cover 14.8\,k survivors; projecting them to everyone was weak (AUC 0.63) and is not used for conclusions.
  \item H4 cannot distinguish bans from voluntary exit; the 2022 volume drop is unexplained.
  \item With millions of observations, p-values are uninformative; the evidence is effect sizes and robustness.
  \item Much of the pipeline was AI-written; items tagged \status{[AI, not yet reviewed]} still need your review.
\end{itemize}

% ------------------------------------------------------------------
\section{Open items and next steps}\label{sec:next}
\subsection{Pending from meeting 2}
\begin{itemize}[label=$\square$]
  \item Study the hypothesis tests, methods and results until you can explain them (use Sections~\ref{sec:choices}, \ref{sec:methods} and \ref{sec:selftest}).
  \item Start reading about signed networks, and the group's papers on how people reply to each other (the latter are not in \file{notes/papers/} yet).
  \item Review the decisions tagged \status{[AI, not yet reviewed]} (C17, C18).
  \item Ask the TFG coordinator about cluster access (deferred: not needed yet).
\end{itemize}

\subsection{Optional directions raised in the notes}
\begin{itemize}
  \item Is a vote-based position a valid proxy for opposition? Is the co-present null the right baseline?
  \item How to handle the circularity of faction hostility (factions built from the same votes).
  \item What causes the 2022 volume drop (collection vs platform).
  \item Separating bans from voluntary exits (heavy-burst accounts).
  \item Next direction: \emph{bridging users} (reply across factions without hostility; do they live longer?), \emph{escalation as a sequence} ($P(\text{hostile at depth }k+1\mid\text{hostile at depth }k)$), or topic-level H2/H3.
\end{itemize}

% ------------------------------------------------------------------
\section{Self-test}\label{sec:selftest}
If you cannot answer one of these in two minutes, go to the section in brackets.
\begin{enumerate}
  \item Why can't ``position'' be measured from text here, and what replaces it? [\ref{sec:data}, C2]
  \item What exactly is a reply edge, and what is lost with that definition? [C3]
  \item Why compare with a co-present null instead of with random users? [C5]
  \item What does stratifying by comment buy you in H1? Why is OR 4.4 not ``4.4 times more likely''? [C6, \ref{sec:methods}]
  \item Why exclude A's vote on the specific comment in H2? [C8]
  \item Why per-edge metrics instead of the pooled downvote share? [C9]
  \item What is right-censoring and what does it do to the top 10\% of lifetimes? [C11, \ref{sec:results}]
  \item What is immortal-time bias and how does a landmark at day 30 avoid it? [C12]
  \item What does ``HR 1.11 per SD'' mean, and what does it not mean? [C13, \ref{sec:conclusions}]
  \item What bug did the first version of the shrinkage have, and how did you notice? [C14]
  \item What does the multiverse show, and what would make you worry? [C15]
  \item Why is faction hostility partly circular? [\ref{sec:conclusions}]
  \item Why are the positive-vote results ``the mirror image'' and not new evidence at the comment level? [\ref{sec:results}]
  \item Why define the long-lived group within cohort? [C17]
\end{enumerate}

% ------------------------------------------------------------------
\section{Reading list}\label{sec:reading}
\textbf{Priority 1} (read first, in this order): [\ref{ref:dataset}], [\ref{ref:faultana}], [\ref{ref:steegen}], [\ref{ref:simonsohn}], [\ref{ref:heider}]/[\ref{ref:cartwright}], then one survival text [\ref{ref:kleinbaum}] and one categorical-data text [\ref{ref:agresti}].

\textbf{Local PDFs in \file{notes/papers/}} (not all checked here): Fraxanet 2024 (+SI), Fraxanet 2025 dataset, Lan Ceqin 2026 thesis, Steegen 2016, Simonsohn 2020 (+SI), Cartwright \& Harary 1956, Heider 1946. Also present: D'Ignazi 2025 (ranking and polarization) and Amarasinghe 2026 (misinformation prevalence and uncertainty); the notes do not link them to this analysis.

\subsection*{This group and this data}
\begin{enumerate}[label={[\arabic*]},ref=\arabic*,series=refs]
  \item\label{ref:dataset} Fraxanet, E., G\'omez, V., Kaltenbrunner, A., Pellert, M. (2026). \emph{A Decade of News Forum Interactions: Threaded Conversations, Signed Votes, and Topical Tags.} Scientific Data. doi:10.1038/s41597-026-07267-3. Dataset: doi:10.82201/RVQGUG.
  \item\label{ref:faultana} Fraxanet, E., Pellert, M., Schweighofer, S., G\'omez, V., Garcia, D. (2024). Unpacking polarization: Antagonism and alignment in signed networks of online interaction. \emph{PNAS Nexus}. (Local PDF; verify volume and pages.)
  \item\label{ref:gomez08} G\'omez, V., Kaltenbrunner, A., L\'opez, V. (2008). Statistical analysis of the social network and discussion threads in Slashdot. \emph{Proc. WWW}.
  \item\label{ref:gomez13} G\'omez, V., Kappen, H. J., Litvak, N., Kaltenbrunner, A. (2013). A likelihood-based framework for the analysis of discussion threads. \emph{World Wide Web} 16.
\end{enumerate}

\subsection*{Signed networks and balance}
\begin{enumerate}[label={[\arabic*]},ref=\arabic*,resume=refs]
  \item\label{ref:heider} Heider, F. (1946). Attitudes and cognitive organization. \emph{Journal of Psychology} 21, 107--112.
  \item\label{ref:cartwright} Cartwright, D., Harary, F. (1956). Structural balance: a generalization of Heider's theory. \emph{Psychological Review} 63(5), 277--293.
  \item\label{ref:leskovec} Leskovec, J., Huttenlocher, D., Kleinberg, J. (2010). Signed networks in social media. \emph{Proc. CHI}; and Predicting positive and negative links in online social networks. \emph{Proc. WWW}.
  \item\label{ref:easley} Easley, D., Kleinberg, J. (2010). \emph{Networks, Crowds, and Markets.} Cambridge University Press. (Chapter on positive and negative relationships; freely available online from the authors.)
  \item\label{ref:wasserman_net} Wasserman, S., Faust, K. (1994). \emph{Social Network Analysis: Methods and Applications.} Cambridge University Press.
\end{enumerate}

\subsection*{Statistics: stratified tables, resampling, Bayesian shrinkage}
\begin{enumerate}[label={[\arabic*]},ref=\arabic*,resume=refs]
  \item\label{ref:mh} Mantel, N., Haenszel, W. (1959). Statistical aspects of the analysis of data from retrospective studies of disease. \emph{J. National Cancer Institute} 22(4), 719--748.
  \item\label{ref:rbg} Robins, J., Breslow, N., Greenland, S. (1986). Estimators of the Mantel--Haenszel variance consistent in both sparse data and large-strata limiting models. \emph{Biometrics} 42, 311--323.
  \item\label{ref:agresti} Agresti, A. \emph{Categorical Data Analysis.} Wiley (3rd ed., 2013).
  \item\label{ref:efron} Efron, B. (1979). Bootstrap methods: another look at the jackknife. \emph{Annals of Statistics} 7(1), 1--26.
  \item\label{ref:bda} Gelman, A., Carlin, J., Stern, H., Dunson, D., Vehtari, A., Rubin, D. \emph{Bayesian Data Analysis} (3rd ed.). CRC Press, 2013. (Beta--binomial model; hierarchical shrinkage.)
  \item\label{ref:wasserstein} Wasserstein, R., Lazar, N. (2016). The ASA statement on p-values: context, process, and purpose. \emph{The American Statistician} 70(2), 129--133.
  \item\label{ref:lin} Lin, M., Lucas, H., Shmueli, G. (2013). Too big to fail: large samples and the p-value problem. \emph{Information Systems Research} 24(4).
  \item\label{ref:hernan} Hern\'an, M., Robins, J. (2020). \emph{Causal Inference: What If.} Chapman \& Hall/CRC. (Free online.)
\end{enumerate}

\subsection*{Survival analysis}
\begin{enumerate}[label={[\arabic*]},ref=\arabic*,resume=refs]
  \item\label{ref:km} Kaplan, E. L., Meier, P. (1958). Nonparametric estimation from incomplete observations. \emph{J. American Statistical Association} 53, 457--481.
  \item\label{ref:cox} Cox, D. R. (1972). Regression models and life-tables. \emph{J. Royal Statistical Society B} 34, 187--220.
  \item\label{ref:kleinbaum} Kleinbaum, D., Klein, M. \emph{Survival Analysis: A Self-Learning Text.} Springer. (Best first read: censoring, KM, Cox, PH assumption.)
  \item\label{ref:schoenfeld} Schoenfeld, D. (1982). Partial residuals for the proportional hazards regression model. \emph{Biometrika} 69, 239--241.
  \item\label{ref:grambsch} Grambsch, P., Therneau, T. (1994). Proportional hazards tests and diagnostics based on weighted residuals. \emph{Biometrika} 81, 515--526.
  \item\label{ref:anderson} Anderson, J., Cain, K., Gelber, R. (1983). Analysis of survival by tumor response. \emph{J. Clinical Oncology} 1, 710--719. (Landmark method.)
  \item\label{ref:suissa} Suissa, S. (2008). Immortal time bias in pharmaco-epidemiology. \emph{American J. Epidemiology} 167, 492--499.
\end{enumerate}

\subsection*{Robustness and researcher degrees of freedom}
\begin{enumerate}[label={[\arabic*]},ref=\arabic*,resume=refs]
  \item\label{ref:steegen} Steegen, S., Tuerlinckx, F., Gelman, A., Vanpaemel, W. (2016). Increasing transparency through a multiverse analysis. \emph{Perspectives on Psychological Science} 11(5), 702--712. (Local PDF.)
  \item\label{ref:simonsohn} Simonsohn, U., Simmons, J., Nelson, L. (2020). Specification curve analysis. \emph{Nature Human Behaviour} 4, 1208--1214. (Local PDF + SI.)
\end{enumerate}

\subsection*{Online communities and user lifecycle}
\begin{enumerate}[label={[\arabic*]},ref=\arabic*,resume=refs]
  \item\label{ref:dnm} Danescu-Niculescu-Mizil, C., West, R., Jurafsky, D., Leskovec, J., Potts, C. (2013). No country for old members: user lifecycle and linguistic change in online communities. \emph{Proc. WWW}.
  \item\label{ref:cheng} Cheng, J., Danescu-Niculescu-Mizil, C., Leskovec, J. (2015). Antisocial behavior in online discussion communities. \emph{Proc. ICWSM}.
  \item\label{ref:kraut} Kraut, R., Resnick, P. (2012). \emph{Building Successful Online Communities: Evidence-Based Social Design.} MIT Press.
\end{enumerate}

\subsection*{Software}
\begin{enumerate}[label={[\arabic*]},ref=\arabic*,resume=refs]
  \item\label{ref:duckdb} Raasveldt, M., M\"uhleisen, H. (2019). DuckDB: an embeddable analytical database. \emph{Proc. SIGMOD (demo)}.
  \item\label{ref:lifelines} Davidson-Pilon, C. (2019). lifelines: survival analysis in Python. \emph{Journal of Open Source Software} 4(40), 1317.
\end{enumerate}

% ------------------------------------------------------------------
\appendix
\section{Where things are}
\begin{tabularx}{\textwidth}{@{}l Y@{}}
\toprule
Path & Content \\
\midrule
\file{INDEX.md} & Project map and how to start a new experiment \\
\file{notes/meetings/} & One note per meeting (2026-09-30, 2026-10-19) \\
\file{notes/dataset/} & README, Codebook, download log, \file{data-treatment.md} (cleaning decisions) \\
\file{notes/method/AI-USE.md} & Assessment of AI use and the learning method \\
\file{notes/experiments/01-reply-opposition-lifetime/} & README, PLAN, NOTES (decision log), RESULTS, EXPLICACION, EXPLICACION-RESUMEN, PREGUNTAS, \file{results/} (CSVs, figures, \file{run.log}) \\
\file{code/experiments/01-reply-opposition-lifetime/} & Scripts 00--09, \file{common.py}, \file{h4.py}, \file{plotstyle.py} \\
\file{data/derstandard/} & Raw data and \file{derived/} Parquet files (outside git) \\
\file{notes/papers/} & PDFs of the readings \\
\bottomrule
\end{tabularx}

\section{Glossary}
\begin{description}[style=nextline,leftmargin=0pt]
\item[Edge (reply edge)] A $\to$ B: A's message is a reply to a comment by B.
\item[Co-present null] Replace B by the author of a random earlier comment in the same article.
\item[Negative / positive edge] A's Beta-posterior approval of B, $E[p]$, below 0.4 / above 0.6 (with a minimum number of votes).
\item[Faction] One of the two FAULTANA groups of stable users.
\item[Hostile reply] A reply where the replier downvoted the comment they answer; \emph{friendly}: upvoted.
\item[Lifetime] Last activity minus first activity (comments and votes).
\item[Right- / left-censored] Last activity on or after 2022-07-01 / first activity before 2013-07-01.
\item[Landmark] Day 30 after the first comment: exposures are measured before it, survival after it.
\item[HR] Hazard ratio; $>1$ means leaving sooner.
\item[Multiverse] Re-running an analysis under all defensible choices to see the spread.
\end{description}

\end{document}
